%! Author = boldizsar.poor
%! Date = 26/10/2026

% Preamble
\documentclass[11pt]{article}

% Packages
\usepackage{amsmath}
\usepackage{stmaryrd}
\usepackage{csquotes}
\usepackage{bookmark}
\usepackage{xcolor}

\usepackage{amsthm,thmtools}
\theoremstyle{plain}
\newtheorem{theorem}{Theorem}[section]
\newtheorem{lemma}[theorem]{Lemma}
\newtheorem{proposition}[theorem]{Proposition}\theoremstyle{definition}
\newtheorem{definition}[theorem]{Definition}
\newtheorem{example}[theorem]{Example}
\newtheorem{remark}[theorem]{Remark}

\usepackage[noabbrev,nameinlink]{cleveref}
\usepackage{todonotes}

\usepackage{macros/tikzfig}
\input{macros/floquet.tikzstyles}
\input{macros/floquet.tikzdefs}

\usepackage[backend=biber]{biblatex}
\addbibresource{macros/AdaptiveSteane.bib}

% Document
\begin{document}

\title{Adaptive Steane-QEC from Non-deterministic Fault-tolerant State-preparation Gadgets}
\author{Boldizsár Poór}

\maketitle

\section{Fault-tolerance}

Here we use the definition of \textcite{gottesmanIntroductionQuantumError2009,gottesmanSurvivingQuantumComputer2026}:
\begin{definition}[FT Error Correction Recovery Property]
    Suppose we have an error correction gadget associated with a QECC with distance $2t + 1$.
    The gadget satisfies the \emph{FT Error Correction Recovery Property}, abbreviated \emph{ECRP}, if, whenever $s \leq t$,
    \[
        \tikzfig{figures/EC-A.tikz}
    \]
\end{definition}
\begin{definition}[FT Error Correction Correctness Property]
    Suppose we have an error correction gadget associated with a QECC with distance $2t + 1$.
    The gadget satisfies the \emph{FT Error Correction Correctness Property}, abbreviated \emph{ECCP}, if, whenever $r + s \leq t$,
    \[
        \tikzfig{figures/EC-B.tikz}
    \]
\end{definition}

\printbibliography

\end{document}